\color{unidarkgreen}{\subsection{Управление выключателем}}
\color{black}

\begin{enumerate}[label=\arabic{section}.\arabic{subsection}.\arabic*,  align=left, labelsep=4pt, leftmargin=0pt, itemindent=57pt]

\item
Функциональный блок управления выключателем (УВ) предназначен для подачи сигнала на включение
силового выключателя. Ввод/вывод функции управления выключателем осуществляется либо посредством
программного ключа \textbf{«УВ Вn»}, либо с использованием внешнего сигнала \textit{«Вывод УВ Вn»}.

\item
Команда на включение выключателя «Включить Вn» формируется:
\begin{itemize}
\item от команды оперативного включения выключателя (сигнал \textit{«РКВ Вn»});
\item при срабатывании АПВ (сигнал \textit{«Ср. АПВ Вn»}).
\end{itemize}

\item
Оперативное включение выключателя осуществляется с контролем наличия разрешающего сигнала от
функционального блока КСН. Максимальное время, в течение которого продолжается ожидание условий
оперативного включения, регулируется уставкой \textbf{«Tожид усл. опер. вкл. Вn»}. По истечении данной выдержки
времени команда оперативного включения сбрасывается. Выдержка времени на возврат \textbf{«Тпрод опер. вкл.
Вn»} задает время продления сигнала оперативного включения.

\item
При необходимости может быть введен подхват команды на включение от токового реле,
установленного в цепи ЭМВ, в этом случае сигнал на включение удерживается в сработанном состоянии до
исчезновения тока через ЭМВ (пока блок-контакты выключателя не разорвут цепь включения). Ввод
осуществляется программным ключом \textbf{«Подхват вкл. Вn от РТ ЭМВ»}.

\item
Предусмотрена возможность блокировки многократных включений на КЗ. Если при наличии сигнала
включения будет протекать ток по электромагнитам отключения (что соответствует включению на КЗ), то
сигнал \textit{«Включить Вn»} заблокируется, пока не пропадут сигналы включения (от АПВ или оперативного
включения). Блокировка вводится программным ключом \textbf{«Блок. от многокр. вкл. на КЗ Вn»}.

\item
При наличии сигнала \textit{«Блок. включения Вn»}, формируемого функциональным блоком контроля силового
выключателя, включение выключателя блокируется.

\begin{figure}[H]
\centering
\includegraphics[width=0.8\textwidth,height=0.8\textheight,keepaspectratio]{\fbpath/761.3391 УВ/img/UV.pdf}
\captionof{figure}{Функциональный блок управления выключателем В1}\label{pic_uv:UV}
\end{figure}

\small
\begin{longtable}{|>{\centering\arraybackslash}m{5.3cm}|>{\centering\arraybackslash}m{4.2cm}|>{\centering\arraybackslash}m{1cm}|>{\centering\arraybackslash}m{3.3cm}|>{\centering\arraybackslash}m{1.8cm}|}
% Первая страница: полный заголовок
\caption{Перечень уставочных параметров УВ\hfill\vspace{-0.5\baselineskip}}\label{uv:tbl1}\\
\hline
\rowcolor{gray!30}
Наименование уставки & Диапазон & Шаг & Значение по умолчанию & Единица измерения \\
\hline
\endfirsthead

% Промежуточные страницы: «Продолжение таблицы»
\caption*{\hspace{3pt}\emph{Продолжение таблицы \ref{uv:tbl1}}\hfill\ }\vspace{-0.5\baselineskip} \\
\hline
\rowcolor{gray!30}
Наименование уставки & Диапазон & Шаг & Значение по умолчанию & Единица измерения \\
\hline
\endhead

%===m>УВ|-
%===m
\end{longtable}
\normalsize	

\end{enumerate}